\documentclass[11pt]{article}

\usepackage[margin=1in]{geometry}
\usepackage{xcolor}
\usepackage{hyperref}
\usepackage{revision-marks}

% Point badges back to the response letter PDF basename.
\setresponseletter{responseLetter}
\showrevmarkstrue

\title{Optimizing Neural Networks for Edge Devices:\\A Comprehensive Study}
\author{San Zhang \and Si Li \and Wu Wang \and Liu Zhao}
\date{}

\begin{document}
\maketitle

\begin{abstract}
We study practical optimization techniques for deploying neural networks on
resource-constrained edge devices. This demo manuscript exists so the response-letter
workflow can be compiled end-to-end out of the box.
\end{abstract}

\section{Introduction}
Deploying large neural networks on edge hardware remains challenging because of
tight memory and latency budgets.
\revise[r1c1a]{Our approach is architecture-agnostic: the same optimization
pipeline applies to CNNs, Transformers, and hybrid models without hand-tuned
architecture-specific rules.}

\section{Method}
Let $\lambda$ denote the regularization coefficient that balances accuracy and
compression.
\revise[r2c1]{Formally, $\lambda > 0$ controls the strength of the complexity
penalty in Equation~(5).}

\section{Experiments}
\begin{figure}[h]
\centering
\fbox{\parbox{0.85\linewidth}{\centering\vspace{1.2cm}
Demo Figure 3 (axis labels enlarged)\\[0.4em]
\rpoint{r1c2}%
\vspace{1.2cm}}}
\caption{Readable axis labels after revision.}
\end{figure}

\revise[r1c1b]{Across ResNet, MobileNet, and a compact Transformer backbone, the
method improves the accuracy--latency Pareto front without per-architecture
retuning.}

\section{Writing fixes}
\revise[r3c1]{The network is trained end-to-end with knowledge distillation.}

\end{document}
